% ===================================================================
%  সারণি নং 3 — শৈশব ও যৌবন।
%
%  Ceci n'est PAS une transcription : c'est une traduction, faite en
%  2026, en bengali d'aujourd'hui. Voir texto/bn/10-sarani-01.tex
%  pour la consigne, qui vaut pour les seize fichiers.
%
%  « জুলাই » n'est pas en gras ici, non plus qu'en ido : le fac-simile
%  l'oublie entre « জুন » et « আগস্ট », et c'est gravuri/korekti.json
%  qui le retablit, colonne par colonne, sur la page de lecture.
% ===================================================================

%%K t03-apar-1 apar
\VUsaut{3.0mm}
\VUtitre{15.0pt}{60}{সারণি নং 3}
\VUsaut{1.6mm}
\VUtitre{11.4pt}{0}{শৈশব ও যৌবন।}
\VUsaut{3.4mm}

%%K t03-apar-2 apar
(সারণি 3 ও 4 এমনভাবে সাজানো যে চারটি \VUgras{পারিবারিক দৃশ্যে}
\VUgras{ঋতু}, \VUgras{বয়স}, \VUgras{পরিবার}, \VUgras{পোশাক} ও
\VUgras{স্বাস্থ্যের} মূল শব্দভান্ডার সামনে এসে যায়।)
\VUsaut{4.0mm}

%%K t03-tit-1 sub
\VUcentre{12.0pt}{0}{\textit{প্রথম দৃশ্য।}}
\VUsaut{3.0mm}

%%K t03-tit-2 sub
\VUpk{10.2pt}{40}{গ্রামে এক দীক্ষাস্নান।}
\VUsaut{2.4mm}

%%K t03-tit-3 sub
\VUpk{10.2pt}{40}{বসন্ত।}
\VUsaut{3.0mm}

%%K t03-c1-01-1 p
1. --- আমরা \VUgras{বসন্তে}, \VUgras{মার্চ}, \VUgras{এপ্রিল} বা
\VUgras{মে} মাসে, \VUgras{সপ্তাহের} কোনো দিনে --- \VUgras{সোমবার},
\VUgras{মঙ্গলবার} বা \VUgras{বুধবার}। সকাল \VUgras{দশটা}।

%%K t03-c1-02-1 p
2. --- \VUgras{আবহাওয়া} সুন্দর, \VUgras{বাতাস} নির্মল। সূর্য ঝলমল করছে।
কয়েকটি ছোট ছোট \VUgras{মেঘ} \textsuperscript{(2)} নীল \VUgras{আকাশে}
\textsuperscript{(1)} উঠছে।

%%K t03-c1-03-1 p
3. --- \VUgras{গাছে} এখনও প্রায় \VUgras{পাতা} নেই। রোগা \VUgras{পপলার}
\textsuperscript{(3)} ও উঁচু \VUgras{প্লেন গাছে} \textsuperscript{(17)} এখনও কেবল
কুঁড়ি।

%%K t03-c1-04-1 p
4. --- \VUgras{আবাবিলরা} \textsuperscript{(4)} ফিরে এসেছে, আর আনন্দে চিৎকার
করতে করতে সেই \VUgras{চূড়ার} \textsuperscript{(7)} চারপাশে চক্কর দিচ্ছে যা
\VUgras{ঘণ্টাঘরের} \textsuperscript{(5)}, আর সেই ঘণ্টাঘর দাঁড়িয়ে আছে গ্রামের
\VUgras{গির্জার} \textsuperscript{(6)} উপরে।

%%K t03-tit-4 sub
\VUsaut{3.4mm}
\VUpk{10.2pt}{40}{গির্জা।}
\VUsaut{3.0mm}

%%K t03-c1-05-1 p
5. --- এই গির্জা খুব সাদাসিধা, তার বড় \VUgras{ফটক} \textsuperscript{(13)} ও
তার বড় \VUgras{ঘড়ি} \textsuperscript{(8)} নিয়ে। \VUgras{ছোট কাঁটা}
\textsuperscript{(9)} X সংখ্যার উপরে: সেটি \VUgras{ঘণ্টা} দেখায়।
\VUgras{বড়টি} \textsuperscript{(10)} XII-র (দুপুর) উপরে; সেটি
\VUgras{মিনিট} দেখায়।

%%K t03-c1-06-1 p
6. --- মিনারে \VUgras{ঘণ্টাগুলি} \textsuperscript{(11)} আনন্দে বাজছে। ছাদের
ডগায় পাথরের একটি ছোট \VUgras{ক্রুশ} \textsuperscript{(12)} দাঁড়িয়ে আছে।

%%K t03-c1-07-1 p
7. --- গির্জার কেবল বড় \VUgras{ফটকই} \textsuperscript{(13)} নয়, পাশে একটি ছোট
দরজাও আছে। সেই দরজা দিয়েই \VUgras{যাজক} \textsuperscript{(15)}, তাঁর
\VUgras{আলখাল্লা} পরে, \VUgras{বস্ত্রাগার} \textsuperscript{(14)} থেকে বেরিয়ে
\VUgras{যাজক-নিবাসের} \textsuperscript{(19)} দিকে যাচ্ছেন।

%%K t03-c1-08-1 p
8. --- তাঁর কাছেই এই যে \VUgras{দুধওয়ালি} \textsuperscript{(24)} তার
\VUgras{গাধা} \textsuperscript{(23)} নিয়ে। সে শহর থেকে ফিরছে, যেখানে
ছেলেমেয়েদের জন্য তার ভালো দুধ পৌঁছে দিয়েছে।

%%K t03-tit-5 sub
\VUsaut{4.0mm}
\VUpk{10.2pt}{40}{ফলের বাগান।}
\VUsaut{3.4mm}

%%K t03-c1-09-1 p
9. --- মশাই আলিবোরঁ (গাধা) এই সকালের দৌড়ে খুব খুশি। সে পাশের
\VUgras{বাগানের} সেই \VUgras{ফুলের} \textsuperscript{(20)} গন্ধে ভরা হাওয়া
আনন্দে টেনে নিচ্ছে, যে ফুল \VUgras{ফলগাছগুলিকে} \textsuperscript{(18)} ঢেকে
রেখেছে। সেই \VUgras{পিচগাছ} \textsuperscript{(18)} ও \VUgras{খুবানিগাছ}
\textsuperscript{(19)} একেবারে গোলাপি ও সাদা, আর ভালো ফসলের আশা জাগায়।

%%K t03-c1-10-1 p
10. --- এদিকে ছোট্ট \VUgras{মৌমাছিরা} \textsuperscript{(22)}, যারা
\VUgras{মৌচাক} \textsuperscript{(21)} থেকে আসে, সেখানে গিয়ে গুনগুন করতে করতে
তাদের মিষ্টি \VUgras{মধু} সংগ্রহ করে।

%%K t03-c1-11-1 p
11. --- কিন্তু এ কী হচ্ছে? গ্রামের ছেলেমেয়েরা দৌড়ে আসছে আর ভয় পাওয়া
\VUgras{রাজহাঁসদের} \textsuperscript{(25)} ছত্রভঙ্গ করছে। এই সেই শোভাযাত্রা যা
মুহুরির ছেলের \VUgras{দীক্ষাস্নানের} পরে গির্জা থেকে বেরোচ্ছে।

%%K t03-c1-12-1 p
12. --- ছোট্ট জেমস, তার \VUgras{দোলনায়} \textsuperscript{(26)}, ঘণ্টার আনন্দিত
ধ্বনিতে জেগে ওঠে, আর তার বোন ম্যাডলিন, যে শোভাযাত্রা দেখতে চায়, মাকে বলে
এসে দোরগোড়াতেই তার চুল বেঁধে দিতে ও তাকে জামা পরিয়ে দিতে।

%%K t03-c1-13-1 p
13. --- মা রাজি হন। তিনি নিজের \VUgras{মাথার কাপড়} \textsuperscript{(28)},
নিজের \VUgras{কাঁচুলি} \textsuperscript{(29)} ও নিজের \VUgras{এপ্রন}
\textsuperscript{(30)} পরে দরজার সামনে একটি ছোট চেয়ারে এসে বসেন। ম্যাডলিনের
গায়ে কেবল তার \VUgras{পেটিকোট} \textsuperscript{(31)} ও তার \VUgras{বডিস}
\textsuperscript{(32)}।

%%K t03-c1-14-1 p
14. --- সে কী খুশি! সে তার প্রিয় ফুলগুলি ভুলে যায় --- তার
\VUgras{কার্নেশন} \textsuperscript{(34)}, যার উপরে \VUgras{প্রজাপতিরা}
\textsuperscript{(35)} বসে, তার \VUgras{টিউলিপ} \textsuperscript{(36)}, তার
\VUgras{উপত্যকার লিলি} \textsuperscript{(37)} যা \VUgras{টবে}
\textsuperscript{(38)} আছে, আর সেই টব রাখা \VUgras{দেয়ালের} \textsuperscript{(33)}
উপরে, এবং তার \VUgras{গোলাপ} \textsuperscript{(39)} যা এত সুন্দর বেড়া গড়ে। সে
শোভাযাত্রার মহিলাদের দামি পোশাক দেখছে।

%%K t03-c1-15-1 p
15. --- গ্রামের ছেলেরা পোশাকের দিকে বিশেষ নজর দেয় না। তারা
\VUgras{মিষ্টি} \textsuperscript{(55)} বা \VUgras{পয়সা} \textsuperscript{(52)}
পেতে সামনে ছোটে ও ধাক্কাধাক্কি করে। তাদের একজন কাঁদছে \textsuperscript{(40)}।

%%K t03-c1-16-1 p
16. --- আর একজন \VUgras{কুকুরের} \textsuperscript{(42)} থাবা মাড়িয়ে দিয়েছে,
সে কুঁইকুঁই করে পালাচ্ছে আর \VUgras{মুরগি} \textsuperscript{(43)} ও তার
\VUgras{ছানাদের} \textsuperscript{(44)} উড়িয়ে দিচ্ছে।

%%K t03-tit-6 sub
\VUsaut{5.0mm}
\VUpk{10.2pt}{40}{দীক্ষাস্নানের শোভাযাত্রা।}
\VUsaut{4.0mm}

%%K t03-c1-17-1 p
17. --- শোভাযাত্রার আগে আগে \VUgras{ধাত্রী} \textsuperscript{(45)} চলেছে। তার
মাথায় লম্বা ফিতেওয়ালা সুন্দর \VUgras{টুপি} \textsuperscript{(46)}। সে
\VUgras{নবজাতককে} \textsuperscript{(47)} কোলে নিয়েছে, যে পুরো \VUgras{লেসে}
\textsuperscript{(48)} মোড়া।

%%K t03-c1-18-1 p
18. --- ধাত্রীর পিছনে \VUgras{ধর্মপিতা} \textsuperscript{(49)} ও
\VUgras{ধর্মমাতা} \textsuperscript{(53)} আসছেন। তাঁরা মিষ্টি ও পয়সা বিলোচ্ছেন।
ধর্মপিতার হাতে \VUgras{দস্তানা} \textsuperscript{(51)} ও মাথায় \VUgras{উঁচু
টুপি} \textsuperscript{(50)}। ধর্মমাতা হাতে একটি সুন্দর \VUgras{তোড়া}
\textsuperscript{(54)} নিয়েছেন।

%%K t03-c1-19-1 p
19. --- তাঁদের পিছনে দেখা যায় শিশুটির \VUgras{কাকা} \textsuperscript{(56)} ও
\VUgras{কাকিমাকে} \textsuperscript{(58)}। কাকিমার মাথায় সুরুচিসম্মত
\VUgras{টুপি} \textsuperscript{(59)}, কাকার গায়ে কালো \VUgras{লম্বা কোট}
\textsuperscript{(57)} ও সাদা ওয়েস্টকোট। তারপর আসে \VUgras{খুড়তুতো ভাই}
\textsuperscript{(60)} ও \VUgras{খুড়তুতো বোন} \textsuperscript{(61)}। সে নবজাতকের
\VUgras{বোন} \textsuperscript{(62)}, ছোট্ট বার্থার হাত ধরে আছে।

%%K t03-c1-20-1 p
20. --- বার্থার অন্য \VUgras{ভাই} \textsuperscript{(63)} পিছনে চলেছে। সে এক
\VUgras{আত্মীয়াকে} \textsuperscript{(64)} হাত দিয়েছে। পরিবারের
\VUgras{বন্ধুরা} \textsuperscript{(65)} তাদের পরে আসছে।

%%K t03-tit-7 sub
\VUsaut{4.6mm}
\VUpk{10.2pt}{40}{দর্শকেরা।}
\VUsaut{3.4mm}

%%K t03-c1-21-1 p
21. --- শোভাযাত্রার পাশে এই যে এক \VUgras{ভিখারি} \textsuperscript{(66)}, তার
\VUgras{ক্রাচের} \textsuperscript{(67)} ভরে। সে ভিক্ষা চাইছে।

%%K t03-c1-22-1 p
22. --- অন্য দিকে খুব \VUgras{ভদ্র} \textsuperscript{(68)} একটি ছেলে। সে টুপি
তুলে চেনা লোকদের \VUgras{« নমস্কার »} জানায়।

%%K t03-c1-23-1 p
23. --- গ্রামের লোকেরা দীক্ষাস্নানের শোভাযাত্রা দেখতে ঘর থেকে বেরিয়ে
এসেছে। দেখা যাচ্ছে \VUgras{সুতির টুপি} পরা এক \VUgras{গ্রাম্য পুরুষ}
\textsuperscript{(69)} আর তার পাশে এক \VUgras{গ্রাম্য নারী} \textsuperscript{(70)}।
তারপর এক \VUgras{বৃদ্ধা} \textsuperscript{(71)}, যে তার \VUgras{লাঠির}
\textsuperscript{(72)} ভরে দাঁড়িয়ে। শেষে \VUgras{চক্কিওয়ালা}
\textsuperscript{(73)} তার চওড়া \VUgras{পশমি টুপিতে} \textsuperscript{(74)}, আর
\VUgras{খড়ম} \textsuperscript{(75)} পরা এক তরুণী গ্রাম্য নারী, যে তার শিশুকে
হাঁটতে শেখাচ্ছে।

%%K t03-c1-24-1 p
24. --- দৃশ্যটি খুব প্রাণবন্ত।
\VUsaut{4.0mm}

%%K t03-tit-8 sub
\VUcentre{12.0pt}{0}{\textit{দ্বিতীয় দৃশ্য।}}
\VUsaut{3.0mm}

%%K t03-tit-9 sub
\VUpk{10.2pt}{40}{সর্বজনীন উৎসব।}
\VUsaut{2.4mm}

%%K t03-tit-10 sub
\VUpk{10.2pt}{40}{জলক্রীড়া ও নৌকাবাইচ।}
\VUsaut{3.0mm}

%%K t03-c2-01-1 p
1. --- \VUgras{গ্রীষ্ম} এসে গেছে। \VUgras{জুন} (জুলাই বা \VUgras{আগস্ট})
মাসের তপ্ত সূর্য নবীনদের স্নান করতে ও নৌকা চালাতে প্রলুব্ধ করে।

%%K t03-c2-02-1 p
2. --- নদীর ডান তীরে দাঁড়িরা জলক্রীড়া ও নৌকাবাইচে নামার জন্য জামাকাপড়
খুলছে।

%%K t03-c2-03-1 p
3. --- তাদের একজন নিজের \VUgras{জুতো} \textsuperscript{(1)} খুলছে; সে তার ফিতে
(বোতাম) খুলছে। তার দুই সঙ্গী আগেই তৈরি। তাদের গায়ে এখন কেবল নিজেদের
\VUgras{গেঞ্জি} \textsuperscript{(2)} ও \VUgras{সাঁতারের জাঙিয়া}
\textsuperscript{(3)}। তারা বন্ধুদের অপেক্ষা করতে করতে গল্প করছে। তারা অন্য
তীরে বসা মেলা ও উৎসবের কথা বলছে।

%%K t03-c2-04-1 p
4. --- তাদের পাশে মাটিতে তাদের সঙ্গীদের \VUgras{কাপড়} পড়ে আছে: এক
\VUgras{জোড়া} \VUgras{বুট} \textsuperscript{(4)}, \VUgras{জুতো}
\textsuperscript{(5)}, \VUgras{মোজা} \textsuperscript{(6)}, \VUgras{পশমের}
\textsuperscript{(7)} ও \VUgras{খড়ের} \textsuperscript{(8)} টুপি এবং একটি
\VUgras{টাই} \textsuperscript{(9)}।

%%K t03-c2-05-1 p
5. --- এক তরুণ নিজের \VUgras{কোট} \textsuperscript{(10)} যত্ন করে ভাঁজ করেছে।
সে সেটি একটি তোয়ালের উপরে রাখছে, যাতে তার পাশের জন একটু পরেই তাদের দুজনের
\VUgras{স্যুট} বেঁধে ফেলবে।

%%K t03-c2-06-1 p
6. --- ষষ্ঠ ছেলেটি পিছিয়ে পড়েছে। তার মাথায় এখনও তার \VUgras{টুপি}
\textsuperscript{(11)}। সে নিজের \VUgras{জামাও} \textsuperscript{(12)} খোলেনি,
নিজের \VUgras{কলারও} \textsuperscript{(13)} নয়, নিজের \VUgras{কাফও}
\textsuperscript{(14)} নয়। সে হাতে নিজের \VUgras{ওয়েস্টকোট}
\textsuperscript{(17)} নিয়েছে এবং এখনও নিজের \VUgras{পাতলুন}
\textsuperscript{(16)} ও নিজের \VUgras{গ্যালিস} \textsuperscript{(15)} পরে আছে। সে
পরের দৌড়ের জন্য নিশ্চয়ই তৈরি হতে পারবে না।

%%K t03-c2-07-1 p
7. --- তরুণদের পাশে \VUgras{কাঁটাগাছে} \textsuperscript{(18)} ঘেরা একটি পায়ে
চলা পথ নদীর তীরে নেমেছে, যেখানে বুড়ো জেমস \VUgras{নলখাগড়ার}
\textsuperscript{(19)} মধ্যে সন্ধ্যায় ছাড়ার \VUgras{আতশবাজি}
\textsuperscript{(20)} তৈরি করছে।

%%K t03-tit-11 sub
\VUsaut{4.4mm}
\VUpk{10.2pt}{40}{দৌড়।}
\VUsaut{3.4mm}

%%K t03-c2-08-1 p
8. --- নদীতে নিজেদের ছাতার ছায়ায় কয়েকজন মহিলা একটি \VUgras{ডিঙিতে}
\textsuperscript{(21)} বেরিয়েছেন। তাঁরা দৌড় দেখছেন, যা খুব রোমাঞ্চকর। প্রতিটি
\VUgras{নৌকায়} \textsuperscript{(22)} চারজন \VUgras{দাঁড়ি} \textsuperscript{(24)}
সমস্ত শক্তিতে টানছে, আর \VUgras{কর্ণধার} \textsuperscript{(26)}
\VUgras{হাল} \textsuperscript{(25)} ধরে নাজুক নৌকাটিকে দিশা দিচ্ছে।
\VUgras{বৈঠা} \textsuperscript{(23)} তালে তালে জলে পড়ছে, আর হালকা নৌকাগুলি
চোখ-ধাঁধানো গতিতে ছুটছে।

%%K t03-c2-09-1 p
9. --- কয়েকজন তরুণ সেতুর থামের পাশে \VUgras{মসৃণ বাঁশের}
\textsuperscript{(28)} পুরস্কার জিততে লেগেছে। তাদের একজন নমনীয় পিছল বাঁশের উপরে
সাবধানে এগিয়ে ডগায় বাঁধা ছোট্ট \VUgras{পতাকা} \textsuperscript{(29)} পর্যন্ত
পৌঁছাতে চাইছে। তার হতভাগ্য \VUgras{প্রতিদ্বন্দ্বী} \textsuperscript{(30)} জলে
পড়ে গেছে। সে সাঁতরে তীরে ফিরছে।

%%K t03-c2-10-1 p
10. --- বড় ছেলেরা \VUgras{জলযুদ্ধ} \textsuperscript{(32)} খেলছে,
\VUgras{ঘাটের} \textsuperscript{(33)} পাশে। এক বিরাট \VUgras{ভিড়}
\textsuperscript{(31)} \VUgras{সেতুর} \textsuperscript{(27)} \VUgras{রেলিঙে} ঝুঁকে
আর \VUgras{তীরে} ঠাসা হয়ে এই সব খেলা দেখছে। তবু কিছু দর্শক ফিরে
\VUgras{চড়ার বাঁশের} \textsuperscript{(45)} খেলা দেখছে, বা
\VUgras{পুরস্কার-বিতরণের} দিকে যাওয়া \VUgras{পৌর শোভাযাত্রা} নিহারছে।

%%K t03-c2-11-1 p
11. --- সবার আগে এই যে কয়েকজন \VUgras{ছোকরা} \textsuperscript{(35)}, যারা
\VUgras{বাজনদারদের} \textsuperscript{(36)} আগে আগে ছুটছে; তারপর আসছেন
\VUgras{পৌরপ্রধান} \textsuperscript{(37)}, তাঁর সহকারীরা ও ছোট শহরের
\VUgras{পৌরসভা}। সবার পিছনে চলছে \VUgras{দমকলকর্মীরা} \textsuperscript{(38)}।

%%K t03-tit-12 sub
\VUsaut{5.0mm}
\VUpk{10.2pt}{40}{মেলা।}
\VUsaut{3.6mm}

%%K t03-c2-12-1 p
12. --- \VUgras{মেলার মাঠে} \textsuperscript{(39)} খুব ভিড়। লোকে নানা রকম
\VUgras{তাঁবু} দেখতে আসে: \VUgras{কাঠের ঘোড়া} \textsuperscript{(40)},
\VUgras{সার্কাস} \textsuperscript{(41)}, যেখানে খেলা শুরুও হয়ে গেছে;
\VUgras{সর্বজনীন নাচ} \textsuperscript{(42)}, যেখানে শহরের তরুণ-তরুণীরা
\VUgras{নাচের} সুখ নিচ্ছে।

%%K t03-c2-13-1 p
13. --- শেষ দিকে দেখা যাচ্ছে \VUgras{ষাঁড়ের আখড়া} \textsuperscript{(44)},
যেখানে বীর স্পেনীয় \VUgras{ষাঁড়-যোদ্ধা} ক্রুদ্ধ \VUgras{ষাঁড়ের}
\textsuperscript{(45)} সঙ্গে লড়ে; তারপর \VUgras{দোলনা-রেল} \textsuperscript{(49)}
ও \VUgras{নিশানার তাঁবু} \textsuperscript{(46)}, যেখানে লোকে \VUgras{বন্দুক}
চালানো ও \VUgras{নিশানায়} মারার অভ্যাস করে।

%%K t03-c2-14-1 p
14. --- কাছেই এই যে \VUgras{মেলার রঙ্গমঞ্চ} \textsuperscript{(47)}, যেখানে
\VUgras{প্রহসন} ও \VUgras{নাটক} অভিনীত হয়। তার একেবারে পাশে
\VUgras{দৈত্য} \textsuperscript{(48)} ও \VUgras{বামনের} তাঁবু। প্রথমজনের
\VUgras{উচ্চতা} দুই মিটার পনেরো; দ্বিতীয়জন কষ্টে একটি শিশুর সমান উঁচু।

%%K t03-c2-15-1 p
15. --- শেষে এই যে বড় \VUgras{চিড়িয়াখানা-গাড়ি} \textsuperscript{(50)}, তার
\VUgras{বন্য জন্তুদের} নিয়ে --- তার \VUgras{বাঘকে} \textsuperscript{(51)}
নিয়ে, যার \VUgras{থাবা} বড় বলবান, তার \VUgras{সিংহকে} \textsuperscript{(52)}
নিয়ে, যার \VUgras{কেশর} ঘন, তার দুই \VUgras{জিরাফ} \textsuperscript{(53)} ও
তার \VUgras{হাতিকে} \textsuperscript{(54)} নিয়ে, যে নিজের \VUgras{শুঁড়}
এত চতুরতায় কাজে লাগায়।

%%K t03-c2-16-1 p
16. --- \VUgras{জন্তু-সাধক} \textsuperscript{(55)}, যিনি এই জন্তুদের শেখান,
তাঁবুর দরজায় দাঁড়িয়ে আছেন।
\VUsaut{6.0mm}
\VUfilet{22mm}
